\section{Introduction}
Monitoring daily activity at home, for example in elderly care, fall-risk
screening or rehabilitation, needs sensing that works in the dark, through
clothing, and without recording images of people. Cameras fail the last
requirement, and wearables depend on users remembering to wear and charge them.
Low-cost single-chip mmWave radars (60--81\,GHz) output a sparse 3-D point cloud
with a Doppler velocity for every point. Identity is not visually recoverable
from such data, which makes radar attractive for privacy-preserving human
activity recognition (HAR) \cite{singh2019radhar,ullmann2023survey}.

Radar point clouds differ sharply from the LiDAR and CAD data that most
point-cloud networks are designed for. A frame contains only about 6--21 returns
(Section~\ref{sec:data}). Many returns are static clutter or multipath rather than
the person. The useful signal lies in how the cloud moves over time rather than
in its instantaneous shape. Existing approaches either voxelise the
cloud \cite{singh2019radhar}, which wastes computation on empty cells, or apply
dense-point backbones such as PointNet++, DGCNN or Point Transformer
\cite{qi2017pointnet2,wang2019dgcnn,zhao2021pt}. Those backbones rely on
farthest-point sampling and $k$-NN search, which are slow on embedded processors.
A second, less visible problem concerns evaluation. Radar HAR samples are short
windows cut with large overlap from long recordings, so random window-level
splits place near-identical samples in training and test. Published accuracies
on the same dataset consequently range from about 34\% to 98\%
\cite{cui2023milipoint,zakka2025makgcn}.

This paper studies privacy-preserving HAR from mmWave point clouds with a focus
on two things: models that are light enough for edge deployment, and
measurements that are honest. Our contributions are:
\begin{itemize}
\item \textbf{STPT-Lite}, a factorised spatio-temporal point transformer.
  Attention runs within short temporal segments and then across segment tokens,
  each token is augmented with an explicit radar motion descriptor, and no FPS or
  $k$-NN is needed. With knowledge distillation, noise-robust training and INT8
  weights it reaches
  \rhOursAcc\% on RadHAR with \oursParamsK\,K parameters, \oursMflops\,MFLOPs,
  \oursIntMB\,MB of storage and \oursLat\,ms single-thread CPU latency.
\item \textbf{Radar-specific pre-processing}: moving-point (Doppler-based) person
  centring and spatial gating. Together they remove the static clutter that makes
  up about half of the returns, along with the scene shortcut that clutter
  creates.
\item \textbf{Leakage-free evaluation} on two public datasets: recording-disjoint
  validation on RadHAR and a subject-disjoint split on MiliPoint, where STPT-Lite
  reaches \mpOursAcc\% top-1 on 49 actions performed by unseen people. The same
  architecture scores \mpRandomAcc\% under a leaky random split, which quantifies how
  much the protocol inflates reported accuracy.
\item A study with \textbf{six re-implemented baselines, three seeds, paired
  significance tests, ablations, MiliPoint$\rightarrow$RadHAR transfer,
  robustness to seven controlled perturbation types, and explainability} (temporal attention,
  occlusion and point saliency).
\end{itemize}
Code, pre-processing and all run logs are released for reproducibility.
